\documentclass[12pt]{article}
\usepackage[utf8]{inputenc}
\usepackage[spanish]{babel}
\usepackage{geometry}
\geometry{a4paper, margin=2.5cm}

\title{Resumen del Proyecto: Aplicación de Gestión Financiera}
\author{}
\date{}

\begin{document}

\maketitle

\section{Acceso y Estructura Principal}
La aplicación utiliza un sistema de \textbf{Autenticación Integrada con Google} (SSO) para el acceso seguro de los usuarios. 

La navegación principal se gestiona a través de un menú lateral izquierdo compuesto por cuatro módulos:
\begin{itemize}
    \item \textbf{Gastos} (Módulo por defecto)
    \item \textbf{Cuotas}
    \item \textbf{Ahorros}
    \item \textbf{Suscripciones}
\end{itemize}

\subsection{Plataformas}
La aplicación mantiene \textbf{dos interfaces distintas} sobre la misma API: una versión de escritorio (menú lateral fijo) y una versión móvil (barra de navegación inferior), con rutas y pantallas propias. No son la misma vista adaptada: cada una tiene sus propios componentes y su nivel de funcionalidad puede diferir. \textbf{El desarrollo prioriza actualmente la versión de escritorio}, que es la que refleja el estado descrito en este documento; las diferencias vigentes de la versión móvil se detallan al final.

\subsection{Gastos como Libro Mayor}
Los cuatro módulos no son compartimentos aislados: \textbf{Cuotas y Suscripciones desembocan automáticamente en Gastos}. Al marcar el pago de una cuota o al confirmar el cobro de una suscripción, el sistema inserta por sí mismo el gasto correspondiente en el historial, etiquetado con su \textit{origen} (\texttt{manual}, \texttt{cuota} o \texttt{suscripcion}) y vinculado al registro que lo generó.

Esta decisión cumple dos funciones. Reduce la carga de registro manual a un solo toque para todo el gasto recurrente y comprometido; y consolida en un único lugar el flujo de caja real del usuario, evitando que el gasto fijo quede invisible para el análisis. La contrapartida es que el análisis debe distinguir el gasto \textbf{discrecional} del \textbf{comprometido}, distinción que el módulo de Gastos resuelve mediante el campo \textit{origen}.

El vínculo es bidireccional en lo que respecta a la corrección de errores: revertir el pago de una cuota o destildar una suscripción \textbf{elimina también el gasto que ese registro había generado}, porque se trata de un error de registro y no de un movimiento real. En cambio, si se elimina el producto o el servicio por completo, los gastos ya generados \textbf{permanecen} en el historial: ese dinero sí salió, y borrarlo falsearía el histórico.

\section{Módulo de Gastos}
Al iniciar sesión, el usuario aterriza en el módulo de "Gastos". Esta sección cuenta con dos pestañas de navegación, abriendo por defecto la vista de \textbf{Análisis}.

Este módulo opera además como \textbf{libro mayor} de la aplicación: recibe tanto los gastos registrados manualmente como los generados de forma automática por los módulos de Cuotas y Suscripciones.

\subsection{Dashboard de Análisis}
La pestaña de análisis está dedicada a la visualización de métricas financieras e integra actualmente dos componentes gráficos:

\begin{enumerate}
    \item \textbf{Evolución del Gasto (Gráfico de Líneas):}
    Proyecta los gastos en una escala temporal mensual. Incluye una herramienta de control financiero que permite al usuario ajustar un umbral o límite de presupuesto para monitorear posibles sobregiros, dibujado como línea de referencia sobre la serie.

    \item \textbf{Distribución del Gasto (Diagrama de Anillo):}
    Implementado como un \textit{donut} interactivo: al seleccionar un segmento, o su entrada en la leyenda, éste se destaca y el centro del anillo muestra su porcentaje y su valor.

    El componente responde a \textbf{dos preguntas distintas}, alternables mediante un conmutador, porque no son equivalentes:
    \begin{itemize}
        \item \textbf{Monto:} cuánto dinero se gastó con cada categoría.
        \item \textbf{Frecuencia:} cuántos movimientos se registraron con cada una.
    \end{itemize}

El gráfico admite además \textbf{dos ejes de agrupación} distintos, ya que un mismo gasto responde a dos preguntas que no conviene mezclar:
    \begin{itemize}
        \item \textbf{Método:} con qué medio se pagó (efectivo, débito, crédito, transferencia).
        \item \textbf{Tipo:} de dónde proviene el gasto (directo, cuotas o suscripciones).
    \end{itemize}

    Un cobro de un servicio cargado a la tarjeta de crédito es \textit{Crédito} en el primer eje y \textit{Suscripciones} en el segundo; agruparlos en una sola dimensión perdería ambas lecturas. Los movimientos cuyo medio de pago no se ha declarado se agrupan bajo una categoría propia y se informan al pie del gráfico.
\end{enumerate}

\section{Módulo de Cuotas}
Este módulo está diseñado para llevar un seguimiento preciso de las compras realizadas a plazo. El enfoque del modelo de datos prioriza el flujo de caja recurrente sobre la deuda histórica: 	extbf{el sistema no almacena el precio original de la compra}, sino únicamente el número de cuotas y el valor de cada una.

Las cifras agregadas que muestra la interfaz (total a pagar, abonado y deuda pendiente) se 	extbf{derivan} de esos dos datos como 	extit{cuotas} $	imes$ 	extit{valor}. Conviene notar que esa multiplicación no equivale al precio de lista del producto, ya que incorpora los intereses del financiamiento: representa el desembolso total comprometido, que es la magnitud relevante para el flujo de caja.

\subsection{Proyección de la Carga Mensual}
El módulo incorpora una proyección de lo que el usuario pagará en cuotas durante los próximos doce meses. Dado que cada producto aporta su cuota únicamente mientras le resten pagos, la serie solo puede descender: \textbf{cada escalón hacia abajo corresponde a una deuda que termina}.

Esta lectura responde la pregunta central del módulo, que el mero listado de productos no resuelve: cuándo se libera capacidad de pago y en qué cuantía. Los indicadores que acompañan al gráfico precisan la \textbf{carga mensual} vigente, la \textbf{deuda pendiente} total, la \textbf{próxima liberación} (mes y monto en que disminuye la carga) y la fecha en que el usuario queda \textbf{libre de cuotas}.

\subsection{Registro de Elementos}
Para dar de alta un nuevo pago en cuotas, el usuario debe registrar la siguiente información del producto:
\begin{itemize}
    \item \textbf{Producto:} Nombre o descripción de la compra.
    \item \textbf{Tienda:} Comercio o entidad donde se adquirió.
    \item \textbf{Cantidad de cuotas:} Número total de pagos acordados.
    \item \textbf{Valor de la cuota:} Monto recurrente a pagar en cada periodo.
    \item \textbf{Medio de pago:} Instrumento con el que se paga el compromiso. Es opcional y se declara una sola vez, puesto que no cambia de una cuota a otra; los gastos que el sistema genera al marcar cada pago lo \textbf{heredan} automáticamente.
\end{itemize}

\subsection{Gestión e Interacción}
La interfaz permite interactuar con cada registro listado. Al presionar sobre un elemento, se despliega un menú de acciones para gestionar el estado de la deuda:
\begin{itemize}
    \item \textbf{Marcar cuota:} Acción rápida que registra un pago efectuado, incrementando el contador de cuotas pagadas y actualizando el progreso.
    \item \textbf{Editar:} Permite la modificación integral del registro. El usuario puede actualizar el producto, la tienda, el total de cuotas, el valor de las mismas y el \textbf{número de cuotas pagadas}. Esta última propiedad es crucial para brindar flexibilidad, permitiendo revertir acciones en caso de presionar "Marcar cuota" por accidente.
    Al reducir el contador de cuotas pagadas, el sistema \textbf{elimina también los gastos que esos marcados habían generado} en el libro mayor, de modo que la corrección no deje movimientos fantasma en el historial. El formulario advierte de ello antes de confirmar. El módulo ofrece además un atajo de \textit{Deshacer última cuota} para el caso más frecuente, que evita abrir el formulario completo.
    \item \textbf{Eliminar:} Borra el registro de la base de datos de forma permanente. Los gastos ya generados por el producto \textbf{se conservan} en el historial, puesto que corresponden a dinero efectivamente desembolsado; únicamente pierden la trazabilidad hacia el producto de origen.
\end{itemize}
\section{Módulo de Ahorros}
Este módulo está orientado al seguimiento del capital acumulado y al monitoreo del progreso hacia diferentes objetivos financieros trazados por el usuario.

\subsection{Registro de Movimientos}
La gestión del ahorro se basa en un historial transaccional. El usuario interactúa principalmente con una lista de \textbf{movimientos}, donde puede registrar dos tipos de operaciones que afectan su saldo disponible:
\begin{itemize}
    \item \textbf{Depósitos:} Aportes de capital destinados a incrementar el fondo de ahorro.
    \item \textbf{Retiros:} Extracciones o disminuciones del fondo.
\end{itemize}

\subsection{Visualización y Metas Financieras}
Una vez que el sistema cuenta con movimientos registrados, se habilita un panel de visualización para analizar la evolución del capital. 

El componente principal es un \textbf{Gráfico de Líneas}, que comparte la base técnica del gráfico de Gastos, trazando el capital acumulado en el tiempo. Sin embargo, su lógica de referencia es distinta: en lugar de alertar sobre un "presupuesto" máximo, el gráfico dibuja \textbf{Metas}.
\begin{itemize}
    \item \textbf{Líneas de Umbral Múltiples:} El gráfico proyecta una línea horizontal por cada meta pendiente, ordenadas de la más cercana a la más lejana y diferenciadas por color, con su nombre y monto en la leyenda. Las metas ya completadas se retiran del gráfico. La escala vertical se ajusta para que la meta más alta quede siempre dentro del área visible, de modo que la brecha resulte legible.
    \item \textbf{Orientación y Brecha Visual:} Esta disposición permite al usuario orientarse rápidamente, visualizando de forma clara e intuitiva la brecha existente (cuánto dinero falta) entre su capital actual (la línea de evolución del ahorro) y cada uno de sus objetivos financieros. La tarjeta de saldo complementa esta lectura indicando de forma explícita el monto restante para alcanzar la meta más próxima.
\end{itemize}

\subsection{Naturaleza de las Metas}
Es importante precisar el alcance del concepto: las metas operan como \textbf{líneas de referencia sobre un saldo único}, no como sobres con dinero asignado. Todos los depósitos alimentan un mismo fondo común, y el progreso de cada meta se mide contra ese saldo total. En consecuencia, varias metas pueden aparecer simultáneamente como alcanzadas aun cuando el capital disponible solo permita cubrir una de ellas. El marcado de una meta como completada es, por tanto, una acción manual del usuario.
\section{Módulo de Suscripciones}
Este módulo está enfocado en la gestión simplificada de gastos recurrentes, permitiendo al usuario tener un panorama claro de sus compromisos financieros fijos.

\subsection{Indicadores Agregados}
Junto al carril, el módulo expone cuatro cifras: el \textbf{compromiso mensual}, el \textbf{costo anualizado} (doce mensualidades más los cobros anuales), el \textbf{acumulado histórico} efectivamente pagado —con el servicio que más ha costado hasta la fecha— y el número de cobros \textbf{por revisar}.

El acumulado merece mención aparte: un cargo mensual reducido no resulta significativo hasta que se suman los meses transcurridos desde el alta, y es precisamente esa cifra la que permite juzgar si el servicio sigue justificando su costo. Un carril mensual, por su propia naturaleza, no la deja ver.

\subsection{Alta de Servicios}
El registro de una nueva suscripción está diseñado para ser rápido y de baja fricción. El usuario únicamente necesita ingresar dos datos clave para dar de alta el registro:
\begin{itemize}
    \item \textbf{Nombre:} Identificador del servicio (ej. Netflix, Gimnasio, Spotify).
    \item \textbf{Monto:} Valor del cobro recurrente.
    \item \textbf{Ciclo y fecha de cobro:} Periodicidad (mensual o anual) y día —y mes, en el caso anual— en que se efectúa el cargo. Si el día no existe en un mes determinado, el cobro se registra el último día disponible.
    \item \textbf{Medio de pago:} Instrumento al que llega el cargo. Opcional, se declara una vez y lo heredan todos los gastos generados.
\end{itemize}

\subsection{Seguimiento Mensual}
La lógica de control de este módulo se basa en un sistema de verificación periódica. Cada mes, la interfaz presenta el listado de suscripciones activas junto a un \textbf{sistema de \textit{checklist}}. El usuario interactúa marcando un \textit{check} manual en cada servicio una vez que constata que el cobro ha sido procesado o pagado en el mes en curso. Esta mecánica previene el solapamiento de pagos, el olvido de cancelaciones y los cobros sorpresa en la tarjeta.

Marcar un servicio registra automáticamente el gasto correspondiente en el libro mayor. Destildarlo se interpreta como la corrección de un error de registro y, por consiguiente, \textbf{elimina también ese gasto} del historial.

El cierre de mes se ejecuta mediante una acción explícita de \textbf{reseteo}, que desmarca la totalidad de los servicios para comenzar un nuevo ciclo de verificación. Cabe señalar que el estado de pago es un indicador único y no lleva asociado el mes al que corresponde: el sistema no conserva, por tanto, un registro de qué meses fueron pagados más allá de los gastos generados en el libro mayor, y el reseteo depende de que el usuario lo ejecute.
\section{Brechas Conocidas}
Esta sección consigna las diferencias vigentes entre lo descrito en este documento y el estado de la implementación, con el fin de mantener la trazabilidad de lo pendiente.

\subsection{Versión Móvil}
El desarrollo prioriza la versión de escritorio, por lo que la versión móvil presenta actualmente las siguientes diferencias:
\begin{itemize}
    \item \textbf{Zoom temporal en Gastos:} la versión móvil incorpora un gráfico de evolución con gestos de \textit{pellizco} y arrastre que conmuta automáticamente entre granularidad mensual y diaria. La versión de escritorio no dispone de esta interacción, cuyo equivalente natural sería el desplazamiento con rueda y el arrastre del cursor.
    \item \textbf{Ahorros sin visualización:} la versión móvil no incluye gráfico alguno en el módulo de Ahorros, por lo que las metas y la brecha de capital no son visualizables en ella.
    \item \textbf{Edición de cuotas:} la funcionalidad de edición descrita en el módulo de Cuotas está disponible únicamente en la versión de escritorio.
    \item \textbf{Duplicación de componentes:} el diagrama de anillo existe en dos implementaciones paralelas, una por plataforma, pendientes de unificación.
\end{itemize}

\subsection{Alcance del Modelo}
\begin{itemize}
    \item \textbf{Ausencia de ingresos:} el sistema registra egresos y ahorro, pero no contempla los ingresos del usuario. En consecuencia, el presupuesto mensual es un valor introducido manualmente y almacenado de forma local en cada dispositivo, sin sincronización entre ellos. La aplicación puede informar cuánto se ha gastado respecto de ese límite autoimpuesto, pero no si el mes cierra en equilibrio.
    \item \textbf{Ausencia de edición en otros módulos:} gastos y suscripciones solo admiten alta y baja; su corrección exige eliminar y volver a crear el registro.
    \item \textbf{Trazabilidad de registros previos:} los gastos generados con anterioridad a la incorporación del vínculo con su origen carecen de dicha referencia, por lo que la reversión automática no alcanza a los pagos marcados antes de ese cambio.
\end{itemize}

\end{document}